% 4.3 小群理论
\section{小群理论}

与 4.2 节一样，设 $\mathbf{G}$ 是带不变子群 $\mathbf{T}$ 的群。本节旨在严格地建立求 $\mathbf{G}$ 的所有表示的方案，即第三章 3.7 节为空间群概述过的同一方案。方案如下：

（i）把 $\mathbf{T}$ 的表示分派成星 $i$。

（ii）从星 $i$ 中任意选一个表示 $D^{i}$。

（iii）确定 $i$ 在 $\mathbf{G}$ 中的小群 $\mathbf{K}^{i}$；记住 $\mathbf{K}^{i} = \sum_{\gamma}r_{\gamma}\mathbf{T}$，其中对 $\gamma$ 的求和取遍由满足 $D^{i}_{\gamma} \equiv D^{i}$ 的元素组成的小陪群 $\overline{\mathbf{K}}^{i}$。

（iv）（原书此处无 (iv)，编号如此。——中译者注）

（v）求 $\mathbf{K}^{i}$ 的表示，并把那些限制到 $\mathbf{T}$ 上给出 $D^{i}$ 的倍数者记为 $\Gamma^{i}_{j}$。它们称为 $\mathbf{K}^{i}$ 的\textit{小表示}（small representations）。

（vi）对每个 $j$ 确定 $\Gamma^{i}_{j} \uparrow \mathbf{G}$。做法是按 4.1 节确立的方法，把这里的 $\mathbf{K}^{i}$ 认同为该节的 $\mathbf{K}_{1}$，把这里的 $\Gamma^{i}_{j}$ 认同为该节的表示 $\Gamma^{j}$。

（vii）对每颗星 $i$ 重复 (ii) 到 (vi)。

上述方案的步骤 (i) 到 (iv) 已详细讨论过。为论证步骤 (v) 的合理性，先证明下面的引理。

\noindent\textbf{引理 4.3.1}\quad 若 $\Gamma^{i}_{j}$ 是 $\mathbf{K}^{i}$ 的表示，且 $\Gamma^{i}_{j} \downarrow \mathbf{T}$ 含有 $D^{i}$，则它只含 $D^{i}$（也许不止一次）。

首先注意，由 Frobenius 互反定理 4.1.5，$\Gamma^{i}_{j} \downarrow \mathbf{T}$ 含 $D^{i}$ 当且仅当 $\mathbf{D}^{i} \uparrow \mathbf{K}^{i}$ 含 $\Gamma^{i}_{j}$。

设 $\mathbf{D}^{i} \uparrow \mathbf{K}^{i} = \sum_{j}\lambda^{i}_{j}\Gamma^{i}_{j} = \Gamma$。定理于是断言：$\Gamma^{i}_{j} \downarrow \mathbf{T}$ 含 $D^{i}$ 恰 $\lambda^{i}_{j}$ 次。但由 $\Gamma$ 的定义，$\Gamma \downarrow \mathbf{T} = (\mathbf{D}^{i} \uparrow \mathbf{K}^{i}) \downarrow \mathbf{T}$ 是 $D^{i}$ 的倍数，因此每个含 $D^{i}$ 的 $\Gamma^{i}_{j} \downarrow \mathbf{T}$ 只能含 $D^{i}$，引理得证。

顺带地，我们还进一步证明了 $(|\mathbf{K}^{i}|/|\mathbf{T}|)D^{i} = \Gamma \downarrow \mathbf{T} = (\sum_{j}\lambda^{i}_{j}\Gamma^{i}_{j}) \downarrow \mathbf{T} = \sum_{j}(\lambda^{i}_{j})^{2}D^{i}$，从而 $\sum_{j}(\lambda^{i}_{j})^{2} = |\mathbf{K}^{i}|/|\mathbf{T}|$。

\noindent\textbf{定理 4.3.1}\quad 表示 $\Gamma^{i}_{j} \uparrow \mathbf{G}$（所有 $i$，所有 $j$）全都不可约；并且这样便得到 $\mathbf{G}$ 的所有表示，且一次也不多、一次也不少。

当然注意，对那些 $\overline{\mathbf{K}}^{i}$ 仅由恒等元组成的星 $i$，有 $\mathbf{K}^{i} = \mathbf{T}$；唯一的 $\Gamma^{i}_{j} = D^{i}$，且 $\mathbf{D}^{i} \uparrow \mathbf{G}$ 不可约，因此定理 4.3.1 把 4.2 节的最终结果作为平凡情形包含在内。

我们先证明 $\Gamma^{i}_{j} \uparrow \mathbf{G}$ 不可约。记住 $\mathbf{G} = \sum_{\alpha}p_{\alpha}\mathbf{K}^{i}$，其中 $p_{\alpha}$ 的特征性质是集合 $\{D^{i}_{\alpha}\}$ 构成 $D^{i}$ 的星。

$\Gamma^{i}_{j} \uparrow \mathbf{G}$ 的不可约性立刻由定理 4.1.4 的推论 (ii) 得到。取 $\mathbf{H}_{\alpha\beta} = \mathbf{T}$ 为 $\mathbf{K}^{i}_{\alpha} \cap \mathbf{K}^{i}_{\beta}$ 对一切 $\alpha \neq \beta$ 的公共子群，该推论断言：若
\begin{equation*}
\sum_{t_{a} \in \mathbf{T}}\chi^{i*}_{\alpha}(t_{a})\chi^{i}_{\beta}(t_{a}) = 0, \quad\text{对一切 } \alpha \neq \beta
\tag{4.3.1}
\end{equation*}
则 $\Gamma^{i}_{j} \uparrow \mathbf{G}$ 不可约。这里 $\chi^{i}_{\alpha}(t_{a})$ 是 $t_{a}$ 在由空间 $\Omega_{\alpha}$ 张成的表示（即可约表示 $\lambda^{i}_{j}D^{i}_{\alpha}$）中的特征标；类似地 $\chi^{i}_{\beta}(t_{a})$ 是 $t_{a}$ 在 $\lambda^{i}_{j}D^{i}_{\beta}$ 中的特征标。由引理 4.3.1，$\Gamma^{i}_{j}$ 限制到 $\mathbf{T}$ 上只含 $D^{i}$，因此 $\lambda^{i}_{j}D^{i}_{\alpha} = \lambda^{i}_{j}D^{i}_{\beta}$（对一切 $\alpha, \beta$），从而式 (4.3.1) 的左端是 $\left(\lambda^{i}_{j}\right)^{2}\sum_{t_{a}}\chi^{*}(t_{a})\chi(t_{a})$，其中 $\chi$ 是 $D^{i}$ 的特征标。由于 $D^{i}$ 不可约，此和等于 $|\mathbf{T}| > 0$；但 $\left(\lambda^{i}_{j}\right)^{2} \geq 1$，矛盾——因此该和为零，$\Gamma^{i}_{j} \uparrow \mathbf{G}$ 不可约。

尚需证明：对每颗星 $i$ 和所有小表示 $j$ 重复诱导过程，我们恰好一次地得到 $\mathbf{G}$ 的所有表示。需要考虑两点：第一，这样做是否会漏掉 $\mathbf{G}$ 的某个表示；第二，是否会重复计入某些表示。为排除第二种可能，先证两个引理。

\noindent\textbf{引理 4.3.2}\quad （限制的传递性）
\begin{equation*}
\left(\Delta^{i}_{j} \downarrow \mathbf{K}^{i}\right) \downarrow \mathbf{T} = \Delta^{i}_{j} \downarrow \mathbf{T}.
\tag{4.3.2}
\end{equation*}
这是因为两端由同一组矩阵组成。

\noindent\textbf{引理 4.3.3}\quad （诱导的传递性）
\begin{equation*}
(\mathbf{D}^{i} \uparrow \mathbf{K}^{i}) \uparrow \mathbf{G} \equiv \mathbf{D}^{i} \uparrow \mathbf{G}
\tag{4.3.3}
\end{equation*}
其中，由于 $\mathbf{D}^{i} \uparrow \mathbf{K}^{i} = \sum_{j}\lambda^{i}_{j}\Gamma^{i}_{j}$，式 (4.3.3) 的左端取为 $\sum_{j}\lambda^{i}_{j}(\Gamma^{i}_{j} \uparrow \mathbf{G}) = \sum_{j}\lambda^{i}_{j}\Delta^{i}_{j}$。

由式 (4.2.15) 之后的正文回顾：张成表示 $\mathbf{D}^{i} \uparrow \mathbf{G}$ 的基由函数 $r_{\alpha}\phi_{p} = \phi_{\alpha p}$（$\alpha = 1$ 到 $|\mathbf{G}|/|\mathbf{T}|$，$p = 1$ 到 $d_{i}$）组成，其中 $\langle \phi_{p} |$ 是 $D^{i}$ 的基。把这些函数按 $\alpha$ 分组，即知 $\mathbf{D}^{i} \uparrow \mathbf{G}$ 限制到 $\mathbf{T}$ 上由 $|\mathbf{G}|/|\mathbf{T}|$ 个块组成，每块都是 $D^{i}$。于是 $\mathbf{D}^{i} \uparrow \mathbf{G}$ 分解成不可约表示的形式必为
\begin{equation*}
\mathbf{D}^{i} \uparrow \mathbf{G} \equiv \sum_{j}\lambda^{i}_{j}\Delta^{i}_{j}.
\tag{4.3.4}
\end{equation*}
但 $\Delta^{i}_{j}$ 在 $\mathbf{G}$ 中不可约。因此由 Frobenius 互反定理，$\Delta^{i}_{j} \downarrow \mathbf{T}$ 恰含 $D^{i}$ $\lambda^{i}_{j}$ 次。

\noindent\textbf{引理 4.3.4}
\begin{equation*}
\Delta^{i}_{j} \downarrow \mathbf{T} = \sum_{\alpha}\lambda^{i}_{j}D^{i}_{\alpha}
\tag{4.3.5}
\end{equation*}
其中集合 $\{D^{i}_{\alpha}\}$ 构成 $D^{i}$ 的星。

上面已证明 $\Delta^{i}_{j} \downarrow \mathbf{T}$ 恰含 $D^{i}$ $\lambda^{i}_{j}$ 次；由于 $\Delta^{i}_{j}$ 是从 $\mathbf{K}^{i}$ 到 $\mathbf{G}$ 的诱导表示，故它必然以相同次数含有 $D^{i}$ 的星的全部成员。

现在很容易证明 $\Gamma^{i}_{j} \uparrow \mathbf{G} = \Delta^{i}_{j}$（所有 $i$，所有 $j$）彼此不等价。首先，设给定 $\Delta^{i_{1}}_{j_{1}}$ 与 $\Delta^{i_{2}}_{j_{2}}$，$i_{1} \neq i_{2}$，则由引理 4.3.4 及不同星的表示互不等价的事实，二者不可能等价。其次，设 $i_{1} = i_{2} = i$ 而 $j_{1} \neq j_{2}$：若 $\Delta^{i}_{j_{1}} \equiv \Delta^{i}_{j_{2}}$，则限制到 $\mathbf{T}$ 后仍等价，从而 $\lambda^{i}_{j_{1}} = \lambda^{i}_{j_{2}} = \lambda$，于是 $\Gamma^{i}_{j_{1}}$ 与 $\Gamma^{i}_{j_{2}}$ 限制后都含 $D^{i}$ 恰 $\lambda$ 次；由 $\sum_{j}(\lambda^{i}_{j})^{2} = |\mathbf{K}^{i}|/|\mathbf{T}|$ 及小表示的完全性可以排除这种情形（详细论证见原文），因此 $\Gamma^{i}_{j_{1}} \neq \Gamma^{i}_{j_{2}}$，等价性不成立。

最后还需证明没有遗漏。设 $\Delta$ 是 $\mathbf{G}$ 的任一不可约表示。由 Frobenius 互反定理，$(\Delta \downarrow \mathbf{T}) \uparrow \mathbf{G}$ 含 $\Delta$；而 $\Delta \downarrow \mathbf{T}$ 必含 $\mathbf{T}$ 的某个不可约表示 $D^{i}$，设其所在星为 $i$、次数为 $\mu$。由引理 4.3.3（传递性）与上面的构造，$\mathbf{D}^{i} \uparrow \mathbf{G}$（它等价于 $(\mathbf{D}^{i} \uparrow \mathbf{K}^{i}) \uparrow \mathbf{G} = \sum_{j}\lambda^{i}_{j}(\Gamma^{i}_{j} \uparrow \mathbf{G})$）在限制到 $\mathbf{T}$ 时含 $D^{i}$ 共 $|\mathbf{G}|/|\mathbf{T}|$ 次；按维数比较可知 $\Delta$ 必与某个 $\Gamma^{i}_{j} \uparrow \mathbf{G}$ 等价。因此没有遗漏。

\begin{equation*}
\sum_{i}\sum_{j}\left(\dim \Delta^{i}_{j}\right)^{2} = \sum_{i}\sum_{j}\left(|\mathbf{G}|/|\mathbf{K}^{i}|\right)^{2}\left(\lambda^{i}_{j}\right)^{2}d^{2}_{i} = \sum_{i}\frac{|\mathbf{G}|^{2}}{|\mathbf{K}^{i}||\mathbf{T}|}d^{2}_{i}
\end{equation*}
\begin{equation*}
= \frac{|\mathbf{G}|}{|\mathbf{T}|}\sum_{i}q_{i}d^{2}_{i} = \frac{|\mathbf{G}|}{|\mathbf{T}|}|\mathbf{T}| = |\mathbf{G}|.
\end{equation*}
定理 4.3.1 证毕。
